% Capítulo 2: Marco Teórico
% Propósito: Presentar las bases teóricas, conceptos clave, antecedentes de investigación y el estado del arte que fundamentan el proyecto.

\section{Marco Teórico}
\label{sec:marco-teorico}

El Marco Teórico reúne los antecedentes, conceptos y fundamentos necesarios para comprender los principales elementos que forman parte de TravelDLS. En este apartado se analizan investigaciones relacionadas con la propuesta, se definen los términos utilizados dentro del proyecto y se presentan las bases teóricas que sustentan la gestión del transporte, el seguimiento y la planificación de rutas.

\subsection{Antecedentes del Proyecto}
\label{sec:antecedentes-proyecto}

Los trabajos revisados en los \textbf{\hyperref[sec:antecedentes]{Antecedentes}} ofrecen referentes para la gestión del transporte, el seguimiento de la carga y la planificación de rutas. A partir de sus aportes y de otras investigaciones relacionadas, se examinan las diferencias de TravelDLS y su contribución propuesta al servicio de carga pesada en Nicaragua.

En la gestión del transporte, el \textbf{\hyperref[ant:aragon]{sistema de administración de fletes de Master Logistics}} organiza la información operativa de una empresa \parencite[pp.~2, 9]{aragon2018implementacion}. TravelDLS retoma esta función para atender a distintas transportistas, manteniendo separados sus registros y vinculándolos con la contratación y el seguimiento de los viajes. Con ello, cada empresa podrá compartir con el cliente el estado de su servicio.

Por su parte, ROUTFALL vincula los registros de clientes, pedidos y costos con la elaboración de rutas para empresas de La Unión, El Salvador \parencite[pp.~52--53, 59--60]{flores2017rutas}. Mientras este trabajo se centra en organizar la distribución, TravelDLS plantea un proceso que comprende la publicación, la contratación y el seguimiento de servicios de carga pesada. La planificación del recorrido formará parte de ese proceso para facilitar la coordinación entre empresa, conductor y cliente.

En cuanto al seguimiento, el \textbf{\hyperref[ant:garcia]{sistema de localización de unidades en León}} permite consultar la ubicación del transporte de pasajeros \parencite[pp.~3, 7]{garcia2022desarrollo}. TravelDLS relacionará esa ubicación con el estado del viaje contratado para que el cliente pueda consultar ambos aspectos del traslado de su carga. El \textbf{\hyperref[ant:pauzuoliene]{estudio de adopción tecnológica en empresas logísticas}} respalda la pertinencia de estas funciones, aunque no demuestra su aceptación en Nicaragua \parencite[pp.~988--989]{pauzuoliene2024smart}.

Los estudios de \textbf{\hyperref[ant:aguilar]{distribución en Tegucigalpa}} y de \textbf{\hyperref[ant:hernandez]{software para Transporte Gonzales}} evalúan herramientas de planificación de rutas en escenarios concretos \parencites[p.~56]{aguilar2022optimizacion}[p.~111]{hernandez2024evaluacion}. TravelDLS seguirá ese criterio al comparar Vecino Más Cercano y 2-opt con Google OR-Tools. Además de evaluar los métodos, integrará esta última herramienta en la gestión de los viajes contratados, sin asumir que alcanzará los ahorros de esos estudios.

En esta misma línea, \textcite[pp.~15--17, 28--31]{martinez2020ruteo} abordan la distribución de bebidas en Quito como un problema de ruteo con capacidad y ventanas de tiempo (CVRPTW), mediante agrupamiento k-medoids y un algoritmo genético. En su evaluación obtuvieron una menor distancia respecto de las estimaciones de la operación estudiada \parencite[p.~37]{martinez2020ruteo}. TravelDLS se diferencia en el problema y los métodos seleccionados: se limita al TSP y a las alternativas antes mencionadas. Por tanto, este antecedente orienta la evaluación de recorridos, pero no constituye una referencia de rendimiento directamente comparable.

Como referencia complementaria, \textcite[pp.~20--24]{alvarez2016plataforma} modelan la distribución de lechuga mediante la asignación de vehículos y la ubicación de centros logísticos. Su plataforma es una red física, mientras que TravelDLS es una aplicación web para gestionar servicios y viajes. Las limitaciones de los datos y del supuesto de transportar un solo producto, reconocidas por los autores, muestran la necesidad de interpretar los resultados según las condiciones de cada modelo.

Entre las soluciones más cercanas a la propuesta, la \textbf{\hyperref[ant:sultan]{plataforma de Sultan y colaboradores}} integra reservas, seguimiento y gestión logística, e incluye conductores que eligen pedidos disponibles para reparto \parencite[pp.~1--4]{sultan2022real}. TravelDLS plantea un esquema distinto para la carga pesada: cada transportista asignará sus vehículos y conductores a los viajes contratados. La diferencia radica en mantener las asignaciones bajo la administración de cada empresa, con acceso del cliente a la información de su servicio, y no en integrar funciones que el antecedente ya reúne.

Asimismo, \textcite{lasluisa2026diseno} utilizan el modelo SaaS multitenant para gestionar flotas y compañías en un sistema de rastreo satelital. TravelDLS comparte este enfoque y lo orienta a conectar la publicación y contratación del servicio con el seguimiento del viaje. El aporte esperado está en la continuidad entre la relación comercial y la operación de transporte, no en el uso de una arquitectura multitenant. El alcance consultado de este antecedente no permite descartar que incluya funciones equivalentes.


Además de las investigaciones relacionadas con la gestión del transporte, el seguimiento de unidades y la planificación de rutas, se consideran estudios que han utilizado tecnologías para desarrollar aplicaciones web con diferentes funcionalidades. Estos trabajos permiten conocer las razones por las que sus autores seleccionaron determinadas herramientas y los aportes que pueden retomarse para el desarrollo de TravelDLS.

En relación con el desarrollo de la plataforma, \textcite{nguyen2022building} utilizó Next.js en una aplicación de comercio electrónico por sus capacidades de renderizado del lado del servidor y generación de páginas estáticas. Asimismo, \textcite{kolding2021developing} incorporaron Next.js en BikeViz para organizar y presentar información geográfica mediante una aplicación web. Ambos proyectos aportan referentes para el desarrollo de la interfaz de TravelDLS, que permitirá publicar servicios de transporte y consultar información relacionada con
las operaciones.

Para el desarrollo del backend, \textcite{farinha2025arkivo} seleccionó AdonisJS en la construcción de un sistema de preservación de arte digital debido a su arquitectura MVC, su compatibilidad con TypeScript y las convenciones que facilitan la organización del código. Esta experiencia constituye un referente para la organización de los módulos de TravelDLS. En cuanto al almacenamiento de información,
\textcite{salunke2024performance} compararon PostgreSQL y MySQL, obteniendo resultados favorables para PostgreSQL en las consultas y operaciones concurrentes evaluadas. Su
investigación aporta criterios relacionados con el rendimiento de las bases de datos y la gestión de los registros de la plataforma.

En cuanto a la representación geográfica, \textcite{kolding2021developing} también seleccionaron Mapbox GL JS para BikeViz debido a sus posibilidades de personalización, la documentación disponible y la facilidad para incorporar funciones a los mapas. Mediante esta herramienta desarrollaron una interfaz que permitía explorar información geográfica a través de diferentes capas. Este antecedente se relaciona con TravelDLS, donde Mapbox se utilizará para representar las ubicaciones de los vehículos y los recorridos asociados con los viajes.

En conjunto, los antecedentes revisados permiten identificar aportes relacionados con la administración del transporte, la contratación de servicios, la asignación de recursos, el seguimiento de viajes y la planificación de recorridos. Asimismo, las investigaciones tecnológicas presentan experiencias sobre la selección de herramientas para el desarrollo de aplicaciones web y la gestión de información. TravelDLS retoma estos aportes y los adapta al transporte terrestre de carga pesada en Nicaragua, mediante una plataforma que relacionará la solicitud del servicio con la administración y el seguimiento de cada viaje hasta su finalización.


\subsection{Definición de Términos}

Para comprender la propuesta de TravelDLS es necesario establecer algunos conceptos que se encuentran directamente relacionados con las operaciones que se pretenden gestionar dentro de la plataforma.

La \hypertarget{concepto-gestion-logistica}{\textbf{gestión logística}} constituye uno de los conceptos centrales del proyecto, ya que una operación de transporte no depende únicamente del desplazamiento de una carga. También requiere organizar información sobre vehículos, conductores, recorridos y estados del servicio. Dentro de TravelDLS, este concepto se entenderá como la coordinación de la información y de las actividades necesarias para desarrollar un servicio de transporte desde su asignación hasta el seguimiento del viaje. En sistemas logísticos desarrollados anteriormente, la integración de clientes, conductores, reservas, envíos y seguimiento ha permitido mantener diferentes componentes de la operación dentro de un mismo entorno \parencite[pp. 1-4]{sultan2022real}. Esta idea se relaciona con TravelDLS, donde se busca que los datos del viaje no se manejen de forma aislada, sino como partes de una misma operación.


En este proceso, la \hypertarget{concepto-orden}{\textbf{orden de transporte o servicio de flete}} representa la solicitud de un servicio para trasladar una carga entre los lugares establecidos por el cliente y la empresa transportista. En este registro se relacionan los datos generales de la carga, su origen y destino, la empresa que prestará el servicio, el precio acordado y el estado de la operación. De esta manera, una misma solicitud puede consultarse desde las distintas funciones del sistema, sin manejar por separado la información de contratación, asignación y seguimiento. La necesidad de organizar los procesos de fletes mediante sistemas de información también ha sido abordada en Nicaragua por \textcite{aragon2018implementacion}.

De igual manera, el \hypertarget{concepto-estados}{\textbf{ciclo de atención y los estados de una orden}} permiten representar las distintas etapas por las que pasa una solicitud desde su registro hasta la conclusión del servicio. Este proceso puede incluir la cotización, la aceptación o el rechazo del precio, la asignación de los recursos necesarios y el desarrollo del traslado. En TravelDLS, los estados ayudarán a identificar la situación actual de cada operación y a diferenciar el avance del servicio de los cambios asociados al pago, ya que una orden entregada y una transacción pagada describen situaciones distintas. La organización de pedidos en sistemas de gestión del transporte se describe en \textcite{sepulveda2023propuesta}; por ello, los estados específicos deberán ajustarse a las reglas del negocio que finalmente se aprueben.

Para ejecutar la orden, la \hypertarget{concepto-asignacion}{\textbf{asignación de recursos de transporte}} consiste en vincular la solicitud con los medios que participarán en su ejecución, principalmente el vehículo y el conductor. Esta relación permite que la empresa identifique quién atenderá la orden y a qué unidad corresponde el traslado, además de consultar la información del recorrido y de la carga desde el mismo servicio. La coordinación de flotas y conductores en el transporte de carga pesada se aborda en \textcite{laura2024tecnologias}. En el diseño actual de TravelDLS se contempla registrar categorías de camión, conductores y vehículos; sin embargo, las reglas específicas sobre disponibilidad, capacidad y compatibilidad deberán establecerse conforme al alcance real del proyecto.

A medida que un viaje se encuentra en desarrollo, cobra importancia conocer qué ocurre con el servicio y conservar información sobre su avance. De esta necesidad surgen los conceptos de \hypertarget{concepto-trazabilidad}{\textbf{trazabilidad}} y \hypertarget{concepto-seguimiento-tiempo-real}{\textbf{seguimiento en tiempo real}}. La trazabilidad permite mantener información relacionada con los estados y acontecimientos de un servicio durante su recorrido, mientras que el seguimiento en tiempo real se concentra en ofrecer la ubicación o la información más reciente disponible. \textcite{pauzuoliene2024smart} señalan que las tecnologías de seguimiento permiten conocer la ubicación y el estado de la carga, además de facilitar la estimación de los tiempos de entrega (p. 988). En TravelDLS, estas funciones estarán relacionadas tanto con las necesidades de control de la empresa transportista como con la información que necesita consultar el cliente durante el traslado.


Junto con la información más reciente del viaje, el \hypertarget{concepto-historial}{\textbf{historial de seguimiento}} permite conservar registros sucesivos para reconstruir la secuencia de ubicaciones o acontecimientos de un envío. Mientras la ubicación actual muestra la posición reportada durante el recorrido, el historial reúne los datos almacenados a lo largo del tiempo. En TravelDLS, esta distinción permitirá que el conductor reporte información vinculada a una orden y que el cliente consulte el progreso de su carga. La frecuencia y oportunidad de las actualizaciones dependerán de la conexión a Internet, la disponibilidad del dispositivo y el mecanismo de registro empleado. Como antecedente nacional de localización de unidades de transporte, se considera el trabajo de \textcite{garcia2022desarrollo}.

El seguimiento de una unidad requiere conocer su posición durante el recorrido, por lo que la \hypertarget{concepto-geolocalizacion}{\textbf{geolocalización}} constituye otro concepto relacionado directamente con el proyecto. Esta permite obtener y representar la ubicación geográfica de un vehículo mediante tecnologías de posicionamiento y servicios de mapas. En Nicaragua, \textcite{garcia2022desarrollo} incorporaron servicios de Google Maps en una aplicación destinada a mostrar en tiempo real la ubicación de unidades de transporte (p. 7). Aunque dicho proyecto corresponde al transporte público, demuestra la aplicación de tecnologías de localización dentro de un sistema de transporte desarrollado en el contexto nacional. En TravelDLS, la geolocalización estará vinculada al seguimiento del viaje y a la consulta de la ubicación de la carga.

Además de conocer dónde se encuentra un vehículo, resulta necesario determinar el recorrido que debe seguir. En este contexto aparecen la \hypertarget{concepto-planificacion-rutas}{\mbox{\textbf{planificación de rutas}}} y la \hypertarget{concepto-optimizacion-rutas}{\mbox{\textbf{optimización de rutas}}}. La planificación consiste en definir y ordenar los puntos que deben visitarse durante un recorrido, mientras que la optimización busca comparar distintas alternativas con el propósito de encontrar recorridos que presenten mejores resultados según criterios como distancia, tiempo o costos. Los antecedentes regionales muestran que los resultados obtenidos mediante herramientas de optimización pueden variar según las condiciones del problema y la solución utilizada. \textcite{aguilar2022optimizacion} reportó mejoras en el escenario analizado mediante la reducción de la distancia recorrida (p. 56), mientras que \textcite{hernandez2024evaluacion} obtuvieron resultados diferentes al evaluar alternativas de software para planificación de rutas (p. 111). Estas diferencias justifican que en TravelDLS no se asuma desde el inicio que un método será superior, sino que se realice una comparación entre las alternativas definidas para el proyecto.


Para comparar recorridos entre varios destinos, se puede utilizar una \hypertarget{concepto-matriz}{\textbf{matriz de distancias}}, una estructura organizada en filas y columnas en la que cada posición representa la distancia entre un punto de origen y otro de destino. Esta representación permite que los métodos de optimización utilicen los mismos datos para calcular y comparar el costo total de distintas rutas. \textcite{tuo2024matrix} emplearon una matriz de distancias dentro de un modelo de distribución logística que incorpora información sobre las condiciones de las vías. Asimismo, cuando el criterio de análisis es la duración del trayecto, la matriz puede representar tiempos de desplazamiento; \textcite{lombard2018matrix} estudiaron un problema de ruteo con tiempos de viaje variables y utilizaron una matriz de distribuciones de probabilidad en lugar de una matriz tradicional de distancias. En TravelDLS, la matriz de distancias o tiempos se contempla como insumo para comparar recorridos con múltiples destinos mediante Vecino Más Cercano, 2-opt y OR-Tools. La forma de obtener estos valores deberá definirse durante la implementación según los datos cartográficos y las condiciones consideradas en el prototipo.

La optimización de recorridos conduce al \hypertarget{concepto-tsp}{\textbf{Problema del Agente Viajero}}, conocido como TSP por sus siglas en inglés. Este problema plantea la necesidad de visitar un conjunto de puntos y regresar al lugar de origen procurando reducir la distancia total recorrida. Dentro de las técnicas que pueden utilizarse para construir una solución se encuentra el \hypertarget{concepto-vecino-mas-cercano}{\textbf{Vecino Más Cercano}}, que selecciona progresivamente el siguiente destino tomando como referencia la menor distancia disponible desde el punto actual. \textcite{contreras2021aplicacion} describen la utilización de este método en situaciones relacionadas con transporte y generación de rutas (p. 11). Para TravelDLS, estos conceptos resultan especialmente importantes debido a que el proyecto contempla el análisis de viajes con múltiples destinos y la comparación de diferentes métodos de optimización.


Otra técnica relacionada con la construcción de recorridos es el \hypertarget{concepto-dos-opt}{\textbf{algoritmo 2-opt}}, un procedimiento de búsqueda local que intenta mejorar una ruta ya construida. Para ello, selecciona dos conexiones del recorrido, las sustituye por una combinación diferente y reorganiza el tramo correspondiente para comprobar si disminuye la distancia total. Cuando encuentra una alternativa de menor costo, conserva el cambio y continúa buscando mejoras hasta alcanzar el criterio de finalización establecido. \textcite{uddin2023twoopt} describen este procedimiento en una propuesta aplicada al Problema del Agente Viajero, mientras que \textcite{dacosta2021twoopt} estudian operadores 2-opt como parte de estrategias de búsqueda local para problemas de rutas. En TravelDLS, 2-opt se considera un procedimiento complementario al Vecino Más Cercano: este último construye una ruta inicial mediante elecciones sucesivas y 2-opt puede utilizarla como punto de partida para mejorar el orden de visita. Los resultados de esta combinación podrán compararse con los obtenidos mediante OR-Tools usando los mismos puntos y criterios de evaluación. No obstante, 2-opt no garantiza encontrar siempre la ruta de menor distancia posible, por lo que su utilidad deberá valorarse mediante las pruebas del proyecto \parencite{dacosta2021twoopt,uddin2023twoopt}.

Estos conceptos permiten comprender los principales procesos que TravelDLS busca integrar. Sin embargo, para explicar de manera más amplia cómo estos elementos pueden organizarse dentro de una plataforma tecnológica, es necesario relacionarlos con enfoques teóricos sobre gestión de transporte, transporte de carga pesada y prestación de software.

\subsection{Bases Teóricas}

Las bases teóricas de TravelDLS permiten explicar los enfoques que sustentan tanto la gestión de las operaciones de transporte como la forma en que la plataforma será ofrecida a diferentes empresas. Uno de los principales referentes son los \hypertarget{concepto-tms}{}Sistemas de Gestión de Transporte (TMS), los cuales integran actividades como la planificación, ejecución y seguimiento de los envíos. \textcite[pp.~10--11]{sepulveda2023propuesta} señala que este tipo de sistemas puede incluir funciones como la gestión de pedidos, planificación de rutas, asignación de cargas, seguimiento de vehículos y facturación, además de facilitar el análisis de indicadores relacionados con costos, tiempos de entrega y utilización de recursos. Estas características guardan relación con TravelDLS, ya que la plataforma busca mantener conectada la información de viajes, vehículos, conductores, cargas y rutas dentro de un mismo entorno.

Esta necesidad de integración se vuelve más importante al considerar las características propias del \hypertarget{concepto-carga-pesada}{}transporte terrestre de carga pesada, donde las operaciones requieren coordinar vehículos, conductores, recorridos y tiempos de entrega. \textcite[p.~67]{laura2024tecnologias} explica que este tipo de transporte presenta desafíos relacionados con la gestión de flotas, la planificación de rutas y la coordinación de conductores, además de restricciones como límites de carga y horarios de conducción. En este contexto, la tecnología sirve como apoyo para centralizar información y automatizar ciertos procesos, lo cual se relaciona directamente con la propuesta de TravelDLS de reunir varias funciones operativas dentro de una misma plataforma.

La forma en que TravelDLS será ofrecido también encuentra sustento en el modelo \hypertarget{concepto-saas}{}Software como Servicio (SaaS). Bajo este enfoque, la aplicación se encuentra disponible a través de Internet y el proveedor se encarga del mantenimiento y las actualizaciones, evitando que cada organización deba administrar una instalación independiente. \textcite[p.~33]{lasluisa2026diseno} explican que este modelo permite ofrecer el software a múltiples usuarios mediante una infraestructura común. Esta característica resulta especialmente adecuada para TravelDLS, ya que la plataforma está pensada para ser utilizada por distintas empresas transportistas y no por una sola organización.

Relacionado con lo anterior, la \hypertarget{concepto-multitenant}{}arquitectura multitenant permite que varias organizaciones utilicen una misma aplicación manteniendo separada la información correspondiente a cada una. \textcite[pp.~34--36]{lasluisa2026diseno} señalan que este enfoque facilita el uso compartido de recursos y el mantenimiento del sistema, aunque también exige mecanismos de seguridad y control de acceso que eviten que un tenant pueda consultar información que no le corresponde. En TravelDLS, esta separación será necesaria para que cada empresa administre sus propios vehículos, conductores, viajes y cargas, mientras que los clientes solo puedan acceder a la información relacionada con los servicios que han contratado.

Esta característica también requiere considerar la \hypertarget{concepto-separacion}{}separación de información operativa, debido a que las empresas transportistas utilizarán una plataforma compartida, pero cada una deberá mantener diferenciados sus registros. En TravelDLS, esta organización se refleja en las relaciones existentes entre empresas, vehículos, conductores y órdenes, permitiendo identificar a qué organización pertenece cada recurso y qué usuarios están autorizados para consultar o modificar sus datos. La importancia de aplicar mecanismos de control de acceso dentro de entornos multitenant se relaciona con lo expuesto por \textcite[pp.~34--36]{lasluisa2026diseno}, quienes destacan la necesidad de proteger la información de las diferentes organizaciones que utilizan una infraestructura común. De esta manera, TravelDLS buscará combinar el uso compartido de la plataforma con una administración diferenciada de los datos operativos de cada empresa.

Además de los enfoques relacionados con la administración y prestación del software, TravelDLS incorpora fundamentos de optimización matemática para apoyar la planificación de recorridos. En este ámbito se considera \hypertarget{concepto-ortools}{}Google OR-Tools, un conjunto de herramientas de software desarrollado para modelar y resolver diferentes problemas de optimización, incluyendo aquellos relacionados con la planificación de rutas. A diferencia del Vecino Más Cercano y 2-opt, que son procedimientos específicos para construir o mejorar recorridos, OR-Tools proporciona componentes que permiten representar el problema, establecer restricciones y ejecutar distintas estrategias de búsqueda. \textcite{muriyatmoko2024ortools} utilizaron esta herramienta con Python para comparar heurísticas y metaheurísticas aplicadas a un problema de ruteo de vehículos con restricciones de capacidad, mientras que \textcite{salsabila2023ortools} la emplearon en un escenario de distribución logística para analizar estrategias como Guided Local Search y Simulated Annealing.

La documentación oficial de Google señala que OR-Tools permite trabajar con diferentes configuraciones de problemas de ruteo y aplicar procedimientos de búsqueda para obtener soluciones de acuerdo con las restricciones establecidas, aunque estas no siempre garantizan alcanzar el resultado globalmente óptimo \parencite{googleortoolsRouting}. Para TravelDLS, la incorporación de esta herramienta está relacionada con el objetivo de desarrollar un módulo que apoye la planificación de rutas con múltiples destinos, utilizando datos de desplazamiento organizados mediante matrices de distancias o tiempos. Su aplicación permitirá comparar los recorridos generados mediante Vecino Más Cercano y mejorados con 2-opt frente a los resultados obtenidos con OR-Tools, manteniendo criterios de evaluación equivalentes. Esta comparación permitirá analizar el comportamiento de las alternativas consideradas antes de valorar su integración y desempeño dentro del prototipo.

En conjunto, estas bases teóricas permiten entender que TravelDLS no se limita a ser una aplicación para registrar viajes, sino que se apoya en principios de gestión de transporte, centralización de procesos y prestación de software mediante una plataforma compartida. Los TMS aportan el enfoque de organización y seguimiento, el transporte de carga pesada define el contexto operativo y el modelo SaaS junto con la arquitectura multitenant explican cómo el sistema puede ser utilizado por diferentes empresas dentro de una misma solución tecnológica. A estos fundamentos se incorpora la optimización de rutas, que permite analizar recorridos con múltiples destinos mediante procedimientos matemáticos y herramientas especializadas.

Con estos fundamentos se establece la base conceptual que orienta la propuesta de TravelDLS. A partir de ella, el siguiente paso consiste en definir el procedimiento metodológico utilizado para seleccionar y analizar la información que sustenta el proyecto, así como organizar las actividades necesarias para continuar con su desarrollo.
